% !TITLE 蝶恋花·庭院深深深几许
% !AUTHOR 欧阳修
% !DYNASTY 两宋
% !GENRE 词
% !ORDER 3

\begin{poem}
庭院深深深几许\textsuperscript{1}？杨柳堆烟\textsuperscript{2}，帘幕无重数。
玉勒雕鞍\textsuperscript{3}游冶处\textsuperscript{4}，楼高不见章台路\textsuperscript{5}。

雨横风狂\textsuperscript{6}三月暮，门掩黄昏，无计留春住。
泪眼问花花不语，乱红\textsuperscript{7}飞过秋千去。
\end{poem}

\begin{annotations}
\notetext{1}{几许：多少。三个“深”字连用，层层递进——庭院到底有多深？深到望不见外面，深到看不见那个人的身影。}
\notetext{2}{杨柳堆烟：杨柳茂密，远远望去像堆着烟雾。烟，指柳絮纷飞或薄雾笼罩下的朦胧感。}
\notetext{3}{玉勒雕鞍：镶玉的马笼头和雕花的马鞍。代指华贵的车马，暗示贵族公子的座驾。}
\notetext{4}{游冶处：游玩冶游的地方，指歌楼妓馆等娱乐场所。}
\notetext{5}{章台路：汉代长安的一条街名，后泛指繁华的游乐场所。楼虽高，也望不见那条路上的他。}
\notetext{6}{雨横风狂：雨下得猛，风刮得急。横（hèng），蛮横、猛烈。三月将尽，风雨交加，春天眼看着就要过去了。}
\notetext{7}{乱红：零乱的落花。风吹花落，花瓣飞过秋千架——春天就这样走了，人也依旧没有回来。}
\end{annotations}

\begin{appreciation}
《蝶恋花·庭院深深深几许》是欧阳修的代表作之一，写一个深闺女子对不归的游子的思念。李清照酷爱这首词，专门用它作首句仿写了一组《临江仙》——能让李清照出手模仿，可见分量。

“庭院深深深几许？”三个“深”字连用，一层比一层深：第一层是院子深，第二层是帘幕多，第三层是人心深得望不见底。“杨柳堆烟，帘幕无重数”——杨柳堆成烟雾，帘幕一重又一重。这个院子不是家，是一座堡垒，把里面的人关住，把外面的人挡住。

“玉勒雕鞍游冶处，楼高不见章台路。”镶玉的马笼头、雕花的马鞍，是那个人的车马。他在外面游玩冶乐，她站在最高的楼上也望不见那条街。下阕，“雨横风狂三月暮，门掩黄昏，无计留春住”——三月将尽，风雨大作，黄昏被关在门外，春天留不住，人也留不住。

“泪眼问花花不语，乱红飞过秋千去。”她流着泪问花，花不说话，零乱的花瓣飞过秋千架，自顾自地走了。为什么要问花？因为没有人可以问。满院子只有秋千、帘幕、落花陪着她。

读到这里，我想起《氓》里的那个女子。她还能爬上墙头眺望复关，看得见的时候载笑载言，看不见的时候哭得痛快——她的喜怒是自己的。而这首词里的女子，连墙都不必爬了：庭院太深，深到她的等待无人知晓，深到春天走了，她都没处说。从《诗经》到北宋，一千多年过去，被关在门里的女人还在等。

好在李清照接住了欧阳修这首词，却没有被庭院关住：她写词、喝酒、划船，出门见识山河。同样是女子，她把“庭院深深”过成了“云中谁寄锦书来”的辽阔。

这首词最难受的地方，不是“他不回来”，是“她出不去”。现在没有人能把你关进庭院了，所以心里堵的时候，别关起门来等一个消息。出去走走，见见朋友，吹吹风，比“泪眼问花”管用。
\end{appreciation}
